%%%%%%%%%%  scan idx 64  |  folio 51  %%%%%%%%%%
\SGAPrintPage{51}%
\sgaStableTarget{sga6.e00.lem.2-9}\textbf{Lemma 2.9.} \emph{Let $\underline L=(L_i)$ be a graded locally free coherent algebraic sheaf, with finitely many nonzero components. In $\underline K(X)$ we have the formula:}
\sgaStableTarget{sga6.e00.eq.2-20}\begin{equation}\tag{2.20}
 \underline E\left((\overset{p}{\otimes}\underline L)^{\mathrm{alt}}\right)=\wedge^p(\underline E(\underline L)).
\end{equation}
In this formula, we consider $\overset{p}{\otimes}\underline L$ as a sheaf on which the group $\mathfrak S_p$ operates, taking the gradings into account (which therefore introduce signs in the well-known way, cf.\ Cartan--Eilenberg); moreover, if $\underline M$ is a graded coherent algebraic sheaf on $X$, we put $\underline E(\underline M)=\sum_{i\geq0}(-1)^i\gamma(M_i)$.

To prove formula (2.20), we are reduced to proving the analogous formula in $\underline K(G)$, where $G$ is an algebraic group (a product of groups $\operatorname{Gl}(n,k)$, if desired) and $\underline L$ a $G$-module of finite dimension over $k$ (cf.\ Chap.\ I, \S 2, example 1, for the definition of $\underline K(G)$). In this case, this formula has already been used in theorem 1.4 to show that the second member of (1.33) satisfies the same functional equation as the first member; its verification is elementary.
\begin{flushright}q.e.d.\end{flushright}

Let $f$ be a morphism from $X$ to $Y$, where $Y$ is an algebraic space. The notion of inverse image of a vector bundle defines a homomorphism $f^!:\underline K(\xi(Y))\longrightarrow\underline K(\xi(X))$ of $\lambda$-rings. Identifying vector bundles with locally free algebraic sheaves, the notion of inverse image of a vector bundle corresponds to the notion of inverse image of an algebraic sheaf (already used in [3, no.\ 8]). When $Y$ is nonsingular, we define more generally a group homomorphism:
\sgaStableTarget{sga6.e00.eq.2-21}\begin{equation}\tag{2.21}
 f^!:\underline K(Y)\longrightarrow\underline K(X)
\end{equation}
%%%%%%%%%%  scan idx 65  |  folio 52  %%%%%%%%%%
\SGAPrintPage{52}%
by a
formula involving an alternating sum of $\Tor$, inspired by the developments at the beginning of \S 2, which the reader will make explicit. The following diagram then commutes:
\[
\begin{tikzcd}[column sep=large,row sep=large]
\underline K(\xi(Y)) \arrow[r,"f^!"] \arrow[d] & \underline K(\xi(X)) \arrow[d]\\
\underline K(Y) \arrow[r,"f^!"] & \underline K(X){,}
\end{tikzcd}
\]
which shows in particular, when $X$ and $Y$ are both nonsingular and quasi-projective, that the two definitions of $f^!$ coincide under the identification of $\underline K(X)$ with $\underline K(\xi(X))$ and of $\underline K(Y)$ with $\underline K(\xi(Y))$.

Another case in which a homomorphism (2.21) can be defined, without $Y$ necessarily being nonsingular, is that in which the functor $F\to f^!(F)$ from $\cF(Y)$ to $\cF(X)$ is exact, i.e.\ when the local rings of $X$ are \emph{flat} modules over the local rings of $Y$. It then suffices to put $f^!(\gamma_Y(F))=\gamma_X(f^!(F))$.
This case occurs in particular when $X$ is a locally trivial fibre space over $Y$ (to see this, we are reduced to the case of a product $X=T\times Y$). Note that we also verify that, in this case, we have $f^!(\underline{O}_Z)=\underline{O}_{f^{-1}(Z)}$ if $Z$ is an irreducible subvariety of $Y$, whence:
\sgaStableTarget{sga6.e00.eq.2-22}\begin{equation}\tag{2.22}
 f^!(\gamma_Y(Z))=\gamma_X(f^{-1}(Z))
\end{equation}
\emph{if $X$ is ``flat'' over $Y$}. Without this last condition, formula (2.22) is in general incorrect, but we have the following analogue of proposition 2.7:

\sgaStableTarget{sga6.e00.prop.2-8}\textbf{Proposition 2.8.}
\emph{Let $f:X\longrightarrow Y$ be a morphism of nonsingular varieties, and $Z$ a cycle on $Y$. If $Z$ is everywhere transverse to $f$, formula}
%%%%%%%%%%  scan idx 66  |  folio 53  %%%%%%%%%%
\SGAPrintPage{53}%
\emph{(2.22) is correct; if $Z$ is homogeneous of codimension $i$ and $f^{-1}(Z)$ has no excess component, we have:}
\sgaStableTarget{sga6.e00.eq.2-23}\begin{equation}\tag{2.23}
 f^!(\gamma_Y(Z))\equiv\gamma_X(f^{-1}(Z)).\qquad\mathrm{mod.}\ \underline K(X)^{i+1}
\end{equation}
Proof analogous to that of proposition 2.7.

Finally, let us mention that we shall see in \S 4 that, \emph{when $X$ and $Y$ are quasi-projective and nonsingular, $f^!$ is compatible with the filtrations of $\underline K(X)$ and of $\underline K(Y)$}.

Again let $X$ and $Y$ be two arbitrary algebraic spaces, and $f$ a \emph{proper} morphism from $X$ to $Y$. For every coherent algebraic sheaf $F$ on $X$, the direct image sheaf $f_!(F)$ is coherent algebraic, and the same holds more generally for the sheaves $R^qf_!(F)$ (cf.\ [3]). Recall that, by definition, $R^qf_!(F)$ is the sheaf on $Y$ associated with the presheaf assigning to an open subset $U$ the group $H^q(f^{-1}(U),F)$; in particular, we have:
\sgaStableTarget{sga6.e00.eq.2-24}\begin{equation}\tag{2.24}
 R^0f_!(F)=f_!(F),\qquad R^qf_!(F)=0\quad(q>\dim X),
\end{equation}
moreover:
\sgaStableTarget{sga6.e00.eq.2-25}\begin{equation}\tag{2.25}
 \Gamma(U,R^qf_!(F))=H^q(f^{-1}(U),F)
\end{equation}
\emph{if the open subset $U$ is affine}. By \S 2, we define a homomorphism from $\underline K(X)$ to $\underline K(Y)$, still denoted $f_!$, by the formula
\sgaStableTarget{sga6.e00.eq.2-26}\begin{equation}\tag{2.26}
 f_!(\gamma_X(F))=\sum_{q\geq0}(-1)^q\gamma_Y(R^qf_!(F)).
\end{equation}
Thus $\underline K(X)$ becomes a \emph{covariant functor in $X$}, with respect to \emph{proper} morphisms. The relation $(fg)_!=f_!g_!$ (just like the relation $(fg)^!=g^!f^!$ that we forgot to mention in its place) is a parti%
%%%%%%%%%%  scan idx 67  |  folio 54  %%%%%%%%%%
\SGAPrintPage{54}%
cular case of the transitivity relations to which we alluded in \S 2, which follow from the spectral sequences of composite functors.

Let us note the following formula, valid for $Y$ nonsingular (and particularly easy to verify when $Y$ is also quasi-projective, by taking for $y$ the class of a vector bundle on $Y$):
\sgaStableTarget{sga6.e00.eq.2-27}\begin{equation}\tag{2.27}
 f_!(x.f^!(y))=f_!(x).y\qquad(x\in\underline K(X),\ y\in\underline K(Y)).
\end{equation}
In the particular case where the inverse images under $f$ of the points of $Y$ are finite and contained in an affine open subset\SGAAuthorFootnote{*}{This last condition is evidently a consequence of the first (something the author still did not know in 1957!).}, we find that the functor $F\longrightarrow f_!(F)$ is exact and that $R^qf_!(F)=0$ for $q>0$; in this case, we therefore have $f_!(\gamma_X(F))=\gamma_Y(f_!(F))$. In particular, if $X$ is an algebraic subspace of $Y$ and $i$ is the canonical injection from $X$ to $Y$, we have:
\sgaStableTarget{sga6.e00.eq.2-28}\begin{equation}\tag{2.28}
 i_!(\gamma_X(F))=\gamma_Y(F),
\end{equation}
where, in the second member, $F$ is considered as a sheaf on $Y$ that is zero outside $X$.

\sgaStableTarget{sga6.e00.sec.4.occ3}\subsection*{\S 4. Some technical results}

Let $X$ be an algebraic space; we saw in \S 3 that the projection $p$ from $X\times k$ onto $X$ defines a homomorphism
\sgaStableTarget{sga6.e00.eq.2-29}\begin{equation}\tag{2.29}
 p^!:\underline K(X)\longrightarrow\underline K(X\times k).
\end{equation}

\sgaStableTarget{sga6.e00.prop.2-9}\textbf{Proposition 2.9.} \emph{The homomorphism (2.29) is surjective.}

We argue by induction on $n=\dim X$; the assertion is trivial if $n<0$; assume $n\geq0$; by proposition 2.6, it suffices to
%%%%%%%%%%  scan idx 68  |  folio 55  %%%%%%%%%%
\SGAPrintPage{55}%
verify that, for every irreducible closed subset $Y$ of $X\times k$, $\gamma(Y)$ is in the image of $p^!$. It suffices to continue the reasoning assuming $X$ irreducible; if $p(Y)$ is not dense in $X$, our assertion follows from the induction hypothesis; otherwise, either $Y=X\times k$ and the assertion is trivial, or $Y$ is a hypersurface in $X\times k$. Now let $U$ be an affine open subset of $X$ consisting of nonsingular points; if $D$ is the divisor on $U\times k$ induced by $Y$, we know (commutative algebra) that it is linearly equivalent to a divisor of the form $D'\times k$, where $D'$ is a divisor on $U$. Then $\gamma_{U\times k}(D)-\gamma_{U\times k}(D'\times k)$ belongs to $\underline K(U\times k)^2$ (corollary of proposition 2.10), and, by what we have already proved, this element is of the form $p^!(x')$ with $x'\in\underline K(U)$. Thus $\gamma_{U\times k}(D)=p^!(y')$ with $y'=x'+\gamma_U(D')$; but $y'$ is the restriction of an element $y$ of $\underline K(X)$ (prop.\ 2.10), and consequently the restriction of $\gamma_{X\times k}(Y)-p^!(y)$ to $U\times k$ is zero; by proposition 2.10, $\gamma_{X\times k}(Y)-p^!(y)$ comes from an element of $\underline K(X\times k-U\times k)=\underline K(X'\times k)$ with $X'=X-U$. The proof is completed by applying the induction hypothesis again.
\begin{flushright}q.e.d.\end{flushright}

\sgaStableTarget{sga6.e00.cor.unnamed.occ4}\textbf{Corollary:} \emph{Let $X$ be an algebraic space and $p$ the projection from $X\times k^n$ onto $X$; then the homomorphism $p^!:\underline K(X)\longrightarrow\underline K(X\times k^n)$ is surjective; it is even bijective when $X$ is nonsingular.}

Induction on $n$ for the first assertion. For the second, we put $f(x)=(x,0)$, whence $pf=\underline{1}$ and $f^!p^!=1$, which proves that $p^!$ is injective.

\sgaStableTarget{sga6.e00.prop.2-10}\textbf{Proposition 2.10.} \emph{Let $X$ be an algebraic space, $U$ an open subset of $X$, and $Y$ its complement in $X$. Then the sequence of natural homomorphisms:}
